The definition rubric achieved 64/120 correct (53.3\%), compared with 69/120
(57.5\%) for the baseline: a decline of 4.2 percentage points. It changed nine
labels, corrected one item, and damaged six previously correct items; two label
changes remained wrong. Stereotype recall decreased from 57.5\% to 52.5\%,
anti-stereotype recall from 15.0\% to 7.5\%, and unrelated recall stayed at 100\%.
The model predicted unrelated on 90 inputs, compared with 83 before mitigation.
Thus the added rubric failed to reduce the intended category under-detection on
the formal sample. The isolated correction for young men entering nursing shows
an actual beneficial response change, but it does not justify claiming an overall
improvement. Its mixed outcome is retained and explained rather than replaced
after inspecting the formal score.

Six-shot alignment achieved \textbf{83/120 correct (69.2\%)}, an improvement of
\textbf{11.7 percentage points} over baseline. Macro-F1 increased from 0.521 to
0.671. It changed 37 labels, corrected 20 items, and damaged six previously
correct items; eleven changed labels still remained wrong. The net gain is
14 correct answers, consistent with $(20-6)/120$. Stereotype recall increased to
87.5\% (35/40), and anti-stereotype recall to 35.0\% (14/40), while unrelated
recall fell to 85.0\% (34/40). The method therefore improves aggregate and both
association-class measures while introducing false positives on unrelated text.
Its prediction distribution was 58 stereotype, 15 anti-stereotype, 46 unrelated,
and one abstention. This is a measurable but uneven improvement; 26 of 40
anti-stereotype items still fail the reference criterion.

Every formal request completed with a single attempt. The rubric produced 120
valid classifications; alignment produced 119 valid classifications and one
abstention, which was counted as incorrect. Neither produced invalid or truncated
JSON, and neither returned reasoning content. Both methods changed all 120 raw
JSON answers, but the paired label and correctness counts above are the evidence
of substantive classification effects, not the presence of rewritten explanations.

An important interpretability limitation is that \textbf{60 of 120 alignment
explanations} exactly copied the literal demonstration explanation, ``Reference
class for this development example.'' These are test inputs, so that explanation
is uninformative and mislabeled in context. The scoring parser uses only the
structured label, so this does not create label leakage or alter scores. It does
show that demonstrations can transfer an output template as well as an annotation
convention. Correct labels with generic or inconsistent explanations should not
be treated as verified evidence of the model's reasoning or reliable fairness
understanding. These responses are preserved verbatim in the example appendix.
